\section{Entropy between two scales}\label{sec:entropy}

We first establish the entropy estimates that turn a separated pair of
translations into a definite gain.  The decisive property of the
self-similar measure is uniform nonsaturation: if its dimension is less
than one, every one-bit scale interval leaves a positive entropy deficit.
The intermediate entropy identities apply to arbitrary compactly
supported probability measures on the line.

For a finite-valued variable $Z$, write
$H(Z)=-\sum_z\P(Z=z)\log\P(Z=z)$, with $0\log0=0$.
Conditional entropy is the average entropy of the conditional law, and
$I(Z;Y)=H(Z)-H(Z\mid Y)$; conditional mutual information is defined
similarly. We use the entropy chain rule and nonnegativity of conditional
mutual information. For a bounded real random variable $X$ and $s>0$, define
\[
 H_s^0(X)=H(\lfloor X/s\rfloor),\qquad
 \mathcal H_s(X)=\int_0^1H(\lfloor X/s+u\rfloor)\,du.
\]
We use the same notation for the law of $X$.  For $s,t>0$, put
\[
 G_{s,t}=\mathcal H_s-\mathcal H_t,
 \qquad \Delta_s=1-G_{s,2s}.
\]
As throughout the paper, all logarithms and entropies are in base~$2$.
Write $h(f)=-\int f\log f$ for differential entropy and $V_s$ for an independent uniform
random variable on $[0,s]$.

The smoothing and nested-grid identities are the averaged-entropy calculus
of \cite[Section~2.2, Lemmas~5, 6, and~10]{VarjuAbsolute2019}; we include their proofs.

\begin{lemma}[Entropy calculus]\label{lem:entropy-calculus}
For every compactly supported real law,
\begin{equation}\label{eq:entropy-smoothing}
 \mathcal H_s(X)=h(X+V_s)-\log s,
 \qquad
 \mathcal H_s(aX+b)=\mathcal H_{s/|a|}(X)
 \quad(a\ne0).
\end{equation}
If $t/s=N$ is a positive integer, then $G_{s,t}$ is concave in the law,
\begin{equation}\label{eq:entropy-integer-window}
 0\le G_{s,t}\le\log N,
 \qquad G_{s,t}(X+Y)\ge G_{s,t}(X)
\end{equation}
for independent bounded $X,Y$.  For arbitrary $z\ge s>0$,
\begin{equation}\label{eq:entropy-scale-comparison}
 \mathcal H_s-\log(z/s+2)\le\mathcal H_z\le\mathcal H_s+1.
\end{equation}
Furthermore,
\begin{equation}\label{eq:entropy-grid-comparison}
 |\mathcal H_s(X)-H_s^0(X)|\le1,
 \qquad |H_s^0(X+b)-H_s^0(X)|\le1.
\end{equation}
\end{lemma}

\begin{proof}
The density of $X+V_s$ is
$f_s(z)=s^{-1}\P(X\in[z-s,z))$ almost everywhere.  It is bounded by
$1/s$ and has bounded support, so its differential entropy is finite,
even when the law of $X$ has atoms.  Set
$p_j(u)=\P((j-u)s\le X<(j+1-u)s)$.  At
$z=(j+1-u)s$ we have $p_j(u)=sf_s(z)$ almost everywhere.
Changing variables in the sum of the integrals
$-p_j(u)\log p_j(u)$ therefore gives
\[
 \mathcal H_s(X)=-\int_{\R}f_s(z)\log(sf_s(z))\,dz.
\]
This proves the first identity.  Translation rotates the grid shift
modulo~$1$, and positive scaling changes the cell length.
Reflection also preserves averaged entropy: $-X+V_s$ is a translate of
$-(X+V_s')$, where $V_s'=s-V_s$ is independent uniform.
These observations prove the signed affine identity.

Suppose $t=Ns$.  Let $A$ be independent uniform on $[0,t)$ and define
\[
 F=\left\lfloor\frac{X+A}{s}\right\rfloor,
 \qquad C=\left\lfloor\frac{X+A}{t}\right\rfloor.
\]
These are jointly translated nested grids, with $C=\lfloor F/N\rfloor$.
Averaging their entropies gives
\begin{equation}\label{eq:nested-conditional-entropy}
 G_{s,t}(X)=H(F\mid C,A).
\end{equation}
Each coarse cell contains $N$ fine cells, proving the two bounds.
For a mixture of laws, conditioning additionally on its mixing variable
can only reduce the right-hand side; this proves concavity, including
for general probability mixtures.  A convolution is a mixture of
translates, so translation invariance proves the convolution inequality.

Finally, for any two translated grids of lengths $s\le z$, a $z$-cell
intersects at most $\lfloor z/s\rfloor+2$ fine cells, and an $s$-cell
intersects at most two coarse cells.  The chain rule bounds the two
conditional entropies by $\log(z/s+2)$ and~$1$.  Averaging gives
\eqref{eq:entropy-scale-comparison}.  The same argument for two grids of
equal length proves \eqref{eq:entropy-grid-comparison}.
\end{proof}

The integer-ratio hypothesis is essential to our use of concavity.
In the argument below it is applied only with ratios $2$, a chosen
integer $K$, or the square of a chosen integer $M$.

\begin{lemma}[Entropy and dimension]\label{lem:entropy-dimension}
Let $\nu$ be a compactly supported exact-dimensional probability measure
on $\R$ of dimension $d$.  Then
\begin{equation}\label{eq:entropy-dimension-upper}
 \limsup_{s\downarrow0}
 \frac{\max\{\mathcal H_s(\nu),H_s^0(\nu)\}}
 {\log(1/s)}\le d.
\end{equation}
\end{lemma}

\begin{proof}
Fix $\eta>0$.  Exact dimensionality gives a measurable set $E$ of mass
at least $1-\eta$ and a common threshold $r_0>0$ such that
\[
 \nu(B(x,r))\ge r^{d+\eta}
 \qquad(x\in E, 0<r<r_0).
\]
Indeed, the sets where this eventual inequality holds below a fixed
reciprocal-integer threshold exhaust a set of full measure; rational
radii and monotonicity suffice to make this construction measurable.

For any translated grid of length $s<2r_0$, choose one point of $E$ in
each cell meeting $E$.  Within each of the three residue classes of
cell indices modulo~$3$, the chosen points have disjoint balls of radius
$s/2$.  Each ball has mass at least $(s/2)^{d+\eta}$.  Thus at most
$3(2/s)^{d+\eta}$ cells meet $E$.  Compact support bounds the total
number of occupied cells by $C/s$, uniformly over grid translations
and small $s$.  Conditioning the cell label on membership in $E$ yields
\[
 H(\text{cell label})
 \le1+\nu(E)\log\bigl(3(2/s)^{d+\eta}\bigr)
       +\nu(E^c)\log(C/s)
 \le(d+2\eta)\log(1/s)+O_{\eta,\nu}(1).
\]
This controls both the fixed grid and its average.  Divide by
$\log(1/s)$ and let $\eta\downarrow0$.
\end{proof}

For Bernoulli convolutions, the following nonsaturation property appears in
Breuillard--Varj\'u \cite[Lemma~13]{BreuillardVarju2019}.
The stopping-time argument below extends it to unequal contractions of
either sign.

\Needspace{5\baselineskip}
\begin{lemma}[Uniform nonsaturation]\label{lem:nonsaturation}
Let $\mu$ be a self-similar probability measure for a finite affine
system $x\mapsto r_i x+t_i$ with $0<|r_i|<1$.  If $\mu$ is exact
dimensional with dimension $d<1$, then
\begin{equation}\label{eq:uniform-nonsaturation}
 \delta:=\inf_{v>0}\Delta_v(\mu)>0.
\end{equation}
No separation or orientation assumption is needed.
\end{lemma}

\begin{proof}
By Lemma~\ref{lem:entropy-calculus}, $0\le\Delta_v\le1$.
We show that a vanishing one-bit deficit would force almost maximal
entropy in a fixed wider scale interval.  Stopped self-similarity then
transfers that estimate to every smaller scale, contradicting the
entropy bound supplied by exact dimensionality.
Suppose there are $v_n>0$ with $\Delta_{v_n}(\mu)\to0$.
For $X\sim\mu$, let $\lambda_v$ be the law of $X+V_v$.
An independent fair bit $J$ gives
\begin{equation}\label{eq:one-bit-information}
 \Delta_v(\mu)
 =h(X+V_{2v})-h(X+V_v)
 =I(J;X+V_v+vJ).
\end{equation}
Thus the equal mixture of $\lambda_v$ and its translate by $v$ carries
little information about its input label.

Here is a quantitative form of that observation.  For two laws $P,Q$,
write $\|P-Q\|_{\mathrm{TV}}=\sup_B|P(B)-Q(B)|$.
Observing whether the mixture sample lies in $B$ bounds the information
in \eqref{eq:one-bit-information} below by
\[
 h_2\bigl((P(B)+Q(B))/2\bigr)
 -\tfrac12h_2(P(B))-\tfrac12h_2(Q(B)),
\]
where $h_2$ is binary entropy.  Since
$h_2''(u)=-1/(\ln(2)u(1-u))\le-4/\ln2$, this expression is at least
$(P(B)-Q(B))^2/(2\ln2)$.  Consequently, for
$e_v=\|\lambda_v-T_v\lambda_v\|_{\mathrm{TV}}$, where $T_v$ denotes
translation by $v$,
\begin{equation}\label{eq:small-information-tv}
 e_v\le\sqrt{2\ln(2)\Delta_v(\mu)}.
\end{equation}

Fix an integer $K\ge2$.  The laws
$P_i=T_{iv}\lambda_v$, $0\le i<K$, have densities $f_i$ and satisfy
$\|P_i-P_0\|_{\mathrm{TV}}\le i e_v$.
Their common density $f_*:=\min_{0\le i<K}f_i$ has mass $m_v$ with
\begin{equation}\label{eq:common-density-mass}
 1-m_v\le\sum_{i=1}^{K-1}\int(f_0-f_i)_+
 \le\frac{K(K-1)}2 e_v.
\end{equation}
For each input label, decompose its law into this same common
submeasure and a residual submeasure of mass $1-m_v$.
The common/residual flag can be sampled independently of the input
label.  On the common branch the output gives no information; on the
residual branch it gives at most $\log K$.  The mutual information of
the uniform mixture is therefore at most $(1-m_v)\log K$.
That mixture is $\mu*\Law(V_{Kv})$, and all its conditional
differential entropies equal $h(X+V_v)$.  For this fixed $K$,
\begin{equation}\label{eq:large-window-saturation}
 G_{v_n,Kv_n}(\mu)\longrightarrow\log K.
\end{equation}

We next transfer this near-saturation to every smaller scale.
Put $a_*:=\min_i|r_i|>0$.  For $0<s<v$, stop the independent coding
at the first prefix whose absolute contraction is at most $s/v$.
The stopping time is bounded because $\max_i|r_i|<1$.
The stopped contractions have absolute values in
$(a_*s/v,s/v]$, and the remaining coding is independent with law $\mu$.
Thus $\mu$ is a finite mixture of affine images of itself with these
contractions.  Signed scaling and integer-ratio concavity give
\begin{align}
 G_{s,Ks}(\mu)
 &\ge\inf_{z\in[v,v/a_*]}G_{z,Kz}(\mu)\notag\\
 &\ge G_{v,Kv}(\mu)-C_* ,
 \qquad C_*:=\log(1/a_*+2)+1.
 \label{eq:stopped-scale-transfer}
\end{align}
For the second inequality, apply
\eqref{eq:entropy-scale-comparison} at $(v,z)$ and $(Kv,Kz)$.
The loss $C_*$ is independent of $K$ and the scales.

Choose an integer $K$ with $(1-d)\log K>C_*+2$.
Then choose one $v=v_n$ such that
$G_{v,Kv}(\mu)>\log K-1$, using
\eqref{eq:large-window-saturation}.  For every $0<s<v$,
\[
 G_{s,Ks}(\mu)>B:=\log K-C_*-1>d\log K.
\]
At $s_m=vK^{-m}$, telescoping gives
\[
 \mathcal H_{s_m}(\mu)-\mathcal H_v(\mu)
 =\sum_{j=1}^mG_{vK^{-j},vK^{-(j-1)}}(\mu)>mB.
\]
Its entropy ratio has lower limit at least $B/\log K>d$, contradicting
Lemma~\ref{lem:entropy-dimension}.  This proves the lemma.
\end{proof}

The next estimate explains how the deficit produces gain.  Unlike the
concavity statements, it does not require an integer ratio of scales.
The fair-pair comparison follows the method of
\cite[Proposition~20 and Lemma~22]{VarjuAbsolute2019}.

\begin{lemma}[Two-point gain]\label{lem:two-point}
Let $X$ be bounded and $J$ an independent fair bit.  For $0<s<v<t$,
\begin{equation}\label{eq:two-point-gain}
 G_{s,t}(X+vJ)-G_{s,t}(X)
 \ge\Delta_v(X)-s/v-v/t.
\end{equation}
The same estimate holds for convolution with any fair two-point law
whose two points have distance $v$.
\end{lemma}

\begin{proof}
Put $Y=X+vJ$.  For each $a>0$, let $U_a$ be independent uniform on
$[0,1)$ and write $Z_a=\lfloor Y/a+U_a\rfloor$.
Translation invariance of the averaged entropy gives
\[
 D_a:=\mathcal H_a(Y)-\mathcal H_a(X)=I(J;Z_a\mid U_a).
\]
Uniform smoothing also gives
$D_v=h(X+V_{2v})-h(X+V_v)=\Delta_v(X)$.

To compare $s$ and $v$, use independent shifts $U_s,U_v$ and the same
sample $Y$.  The information chain rule gives
\[
 I(J;Z_v\mid U_s,U_v)
 \le I(J;Z_s\mid U_s,U_v)
       +H(Z_v\mid Z_s,U_s,U_v).
\]
The first two terms are $D_v$ and $D_s$ because the unused shift is
independent.  For each fixed fine grid and each of its cells of positive
probability, the conditional sample law is independent of $U_v$.
That cell, of length $s<v$, contains a boundary of the $v$-grid on a
proportion $s/v$ of its shifts.  Otherwise the coarse label is fixed;
when a boundary occurs there are at most two coarse labels.  Averaging
therefore gives
$H(Z_v\mid Z_s,U_s,U_v)\le s/v$, whence $D_s\ge D_v-s/v$.

For the coarse-scale gain, adjoining $X$ to the observation yields
\[
 D_t\le I(J;Z_t,X\mid U_t)=I(J;Z_t\mid X,U_t).
\]
For fixed $X=x$, the two labels at $x$ and $x+v$ differ on a proportion
$v/t$ of shifts.  Their conditional information is then one bit and is
otherwise zero.  Thus $D_t\le v/t$.  Subtracting proves
\eqref{eq:two-point-gain}.  The proof allows atoms, since boundary
coincidences occur on null sets of shifts.  Finally, any fair pair is a
translate of $\{0,v\}$ after ordering its two points; translation
invariance proves the last assertion, including for reflected pairs.
\end{proof}
